\section{Descent of the whole fiber}
\label{ds:section}

Let $b=\overline{\C(T_0)}\subset\Omega$ be the algebraically closed
field from Section~\ref{sec:marking}.  Incorporate the finite parameter
change of Proposition~\ref{mk:product} into $I$, and put
\begin{equation}
 L_0=b(I),\qquad k=b(V'),\qquad K=L_0(V')=k(I).
 \label{ds:fields}
\end{equation}
The varieties $I$ and $V'$ are geometrically integral over $b$;
$V'$ is normal and projective.  In particular, $K/k$ is regular.
The generic $x$-fiber $J/K$ is smooth, projective, and geometrically
integral.  Proposition~\ref{mk:product} supplies a smooth projective
geometrically integral reference variety $J_0/k$, with
\begin{equation}
 J_{\overline K}\ \text{is birational to}\ (J_0)_{\overline K},
 \qquad
 \kappa((J_0)_{\overline k})=0,
 \qquad
 h^0((J_0)_{\overline k},pK_{(J_0)_{\overline k}})=1,
 \label{ds:constant-fiber}
\end{equation}
for the ordinary pluricanonical index $p$ used in the rank-one
construction.  Its one-dimensional degree-$p$ space has a nonzero
generator over $k$.

We will descend $K(J)$, after extending $L_0$ to its algebraic closure,
to a field over $b$ that still contains $k=b(V')$.  Thus the coordinates
of the relative Iitaka base remain part of the descended fiber.
The proof has three steps: make the birational descent cocycle constant,
use Lemma~\ref{mk:moving} to remove ramification at moving divisors,
and descend the resulting finite \emph{\'etale} cover.

Here constancy means independence of the parameter variety $I$.
The reference variety $J_0$ is defined over $k=b(V')$, and the
constant cocycle will be defined over a finite extension of $k$.
This dependence on $V'$ is retained when the resulting invariant
field is descended to $b$.

\subsection{Constant birational cocycles}

We use birational groups with the group structure of their
function-field automorphisms.  More explicitly, a birational
selfmap $u$ is identified with $(u^{-1})^*$; products of the resulting
automorphisms are composition.  This convention fixes the order in
the semilinear actions below.

\begin{lemma}[Finite splitting with a constant cocycle]
\label{ds:cocycle}
There are a finite Galois extension $k'/k$ in $\overline k$, a finite
Galois extension $\widetilde K/K$ containing $k'$, and a birational
identification
\[
 J_{\widetilde K}\dashrightarrow (J_0)_{\widetilde K}
\]
with the following properties.  If $G=\operatorname{Gal}(\widetilde K/K)$,
the transported descent action on $\widetilde K(J_0)$ has the form
\begin{equation}
 \rho_g=z_g s_g,
 \qquad
 z_g\in\operatorname{Aut}_{k'} k'(J_0),
 \label{ds:semilinear}
\end{equation}
where $s_g$ transports coefficients by $g$.  The map
\begin{equation}
 G\longrightarrow
 \operatorname{Aut}_{k'} k'(J_0)
       \rtimes\operatorname{Gal}(k'/k),
 \qquad
 g\longmapsto(z_g,g|_{k'})
 \label{ds:faithful-pairs}
\end{equation}
is injective.  In particular,
\begin{equation}
 K(J)\simeq \widetilde K(J_0)^G
 \label{ds:split-field}
\end{equation}
as extensions of $K$.
\end{lemma}

\begin{proof}
\smallskip\noindent\textbf{The birational group.}\enspace
Choose compatible algebraic closures
$\overline k\subset\overline K$.  A nonzero ordinary
pluricanonical form on $J_0$ excludes uniruledness: if a smooth
projective model of $J_0$ were dominated generically finitely by a
ruled variety, its pullback would be a nonzero pluricanonical form
on a resolution of that ruled variety, which has none.

Hanamura's structure theorem for non-uniruled projective varieties
in characteristic zero provides a suitable smooth projective
birational model $T$ of $(J_0)_{\overline k}$ for which the reduced
scheme of flat graphs of birational maps,
\[
 \mathcal B=\Bir(T)_{\mathrm{red}},
\]
is a group scheme locally of finite type over $\overline k$.
Its identity component
\[
 A_1=\mathcal B^0=\operatorname{Aut}^0(T)
\]
is an abelian variety.  This group structure is intrinsic among
such birational models, and represents their algebraic birational
group actions.  These are exactly the conclusions of
\cite[Theorems~2.1--2.2]{Hanamura88}, also stated in
\cite[Theorems~3.8--3.9]{Blanc17}.  The theorem applies to a
non-uniruled projective variety over an algebraically closed
characteristic-zero field; it does not require a minimal model of
$J_0$.

The scheme of flat graphs is an open subscheme of the Hilbert scheme
of $T\times T$; see \cite[Definitions~1.6 and~1.8, Proposition~1.7]{Hanamura87}
and \cite[Proposition~3.7]{Blanc17}.  The Hilbert functor commutes with field extension, and the open graph
condition is geometric. Every birational selfmap over the extended
field is therefore represented by a point of the base-changed graph
scheme. In characteristic zero a reduced locally finite type scheme
is geometrically reduced, so passing to its reduction introduces no
new field-valued points after extension.
Every connected component of $\mathcal B$ contains a
$\overline k$-point, since $\mathcal B$ is locally of finite type
over the algebraically closed field $\overline k$.  Translation by
such a point identifies that component with $A_1$.  Each of these
abelian translates remains connected after extension to
$\overline K$, and the disjoint union of the original open
components exhausts the base-changed scheme.  Thus, after
identifying the birational groups of $T$ and $J_0$, we have
\begin{equation}
 \begin{split}
 \Bir((J_0)_{\overline K})
    &=A_1(\overline K)\Bir((J_0)_{\overline k}),\\
 A_1(\overline K)\cap\Bir((J_0)_{\overline k})
    &=A_1(\overline k).
 \end{split}
 \label{ds:components}
\end{equation}
This argument allows infinitely many components.  The intrinsic
group structure also shows that coefficient transport and
conjugation preserve $A_1$.

\smallskip\noindent\textbf{Removing the nonconstant abelian part.}\enspace
Set $\Gamma=\operatorname{Gal}(\overline K/K)$, and choose an
isomorphism of $\overline K$-function fields
\[
 \theta:\overline K(J)\xrightarrow{\sim}\overline K(J_0).
\]
Let $s_\sigma^J$ and $s_\sigma^0$ denote the natural coefficient
transports on these two fields.  The automorphisms
\[
 u_\sigma=\theta s_\sigma^J\theta^{-1}(s_\sigma^0)^{-1}
       \in\Bir((J_0)_{\overline K})
\]
satisfy
\begin{equation}
 u_{\sigma\tau}=u_\sigma\,\sigma(u_\tau),
 \qquad
 \sigma(v)=s_\sigma^0v(s_\sigma^0)^{-1}.
 \label{ds:cocycle-law}
\end{equation}
The map $\sigma\mapsto u_\sigma$ is continuous when the algebraic
points have the discrete topology: the rational maps defining
$\theta$ and its inverse use only finitely many algebraic
coefficients over $K$.  Its image is therefore finite.

For every component occurring among these values choose a
representative over $\overline k$, choosing the identity in the
identity component.  By \eqref{ds:components} this gives a locally
constant choice $h_\sigma$ and expressions
\[
 u_\sigma=a_\sigma h_\sigma,
 \qquad a_\sigma\in A_1(\overline K).
\]
The component cocycle law and the second equality in
\eqref{ds:components} imply
\[
 d_{\sigma,\tau}
   =h_\sigma\sigma(h_\tau)h_{\sigma\tau}^{-1}
   \in A_1(\overline k).
\]
Consequently
\[
 \alpha_\sigma(a)=h_\sigma\sigma(a)h_\sigma^{-1}
\]
defines an action on $A_1(\overline K)$ preserving
$A_1(\overline k)$.  Indeed the discrepancy between
$\alpha_\sigma\alpha_\tau$ and $\alpha_{\sigma\tau}$ is
conjugation by $d_{\sigma,\tau}$, which acts trivially on the
commutative group $A_1$.

Give
\[
 M_1=A_1(\overline K)/A_1(\overline k)
\]
the discrete topology.  The induced action is continuous: every
algebraic point has an open stabilizer for coefficient transport,
and $h_\sigma$ is locally constant.  In additive notation the
classes $c_\sigma=[a_\sigma]$ satisfy
\begin{equation}
 c_{\sigma\tau}=c_\sigma+\alpha_\sigma(c_\tau).
 \label{ds:additive-cocycle}
\end{equation}

The abelian group $M_1$ is uniquely divisible.  Multiplication by
an integer $n>0$ is surjective on the points of an abelian variety
over either algebraically closed field.  Moreover,
\[
 A_1[n](\overline K)=A_1[n](\overline k),
\]
since the finite torsion scheme is already defined over the
algebraically closed smaller field.  If $na$ is constant, choose
$a_0\in A_1(\overline k)$ with $na_0=na$; then $a-a_0$ is constant
torsion, so $a$ is constant.  This proves injectivity of
multiplication on the quotient as well as its surjectivity.
Thus $M_1$ is a $\Q$-vector space.

Choose an open normal subgroup $N\subset\Gamma$ contained in the
zero set of the continuous cocycle $c$.  Equation
\eqref{ds:additive-cocycle} gives $c_{\sigma n}=c_\sigma$ for
$n\in N$.  It also gives
\[
 \alpha_n(c_\sigma)
 =c_{n\sigma}
 =c_{\sigma(\sigma^{-1}n\sigma)}
 =c_\sigma.
\]
Thus $N$ fixes every cocycle value.  The finite sum
\[
 \beta=\frac{1}{[\Gamma:N]}
            \sum_{\tau\in\Gamma/N}c_\tau
\]
is well-defined, and summing \eqref{ds:additive-cocycle} over the
cosets yields
\begin{equation}
 c_\sigma=\beta-\alpha_\sigma(\beta).
 \label{ds:average}
\end{equation}
Choose $t\in A_1(\overline K)$ lifting $\beta$ and replace $\theta$
by $\theta'=t^{-1}\theta$.  Its cocycle is
\[
 z_\sigma=t^{-1}u_\sigma\sigma(t).
\]
The class of its $A_1$ factor is
$-\beta+c_\sigma+\alpha_\sigma(\beta)=0$ by
\eqref{ds:average}.  Hence
\[
 z_\sigma\in\Bir((J_0)_{\overline k})
 \quad\text{for every }\sigma\in\Gamma.
\]

\smallskip\noindent\textbf{A faithful finite splitting group.}\enspace
The modified cocycle is still continuous, because the algebraic
point $t$ has an open stabilizer.  Its image is therefore finite.
Choose a common finite Galois extension $k'/k$ defining all these
finitely many birational maps and their inverses.  The map
\[
 \Gamma\longrightarrow
 \operatorname{Aut}_{k'} k'(J_0)
          \rtimes\operatorname{Gal}(k'/k),
 \qquad
 \sigma\longmapsto(z_\sigma,\sigma|_{k'})
\]
is a homomorphism, with the multiplication
$(z,\gamma)(w,\delta)=(z\gamma(w),\gamma\delta)$.
It has finite image and an open normal kernel $N_0$.
Put $\widetilde K=\overline K^{N_0}$.  Then
$\widetilde K/K$ is finite Galois and contains $k'$.  For
$\sigma\in N_0$ we have $z_\sigma=1$, so $\theta'$ intertwines the
natural coefficient actions of $N_0$.  Its rational maps descend
to $\widetilde K$ by Galois descent.  The quotient group
$G=\Gamma/N_0$ acts by \eqref{ds:semilinear}, the pairs are
faithful by construction, and the fixed-field identity
\eqref{ds:split-field} follows.
\end{proof}

\subsection{Ramification at moving divisors}

For a finite extension $E/L_0$, a \emph{moving divisor} is a prime
divisor of $V'_E$ whose image is dense in $V'$, as in
Definition~\ref{mk:moving-definition}. Its valuation is trivial
on both $k$ and $E$. We recall the precise conclusion of
Lemma~\ref{mk:moving}.  For the rational degree-$p$ relative
pluricanonical generator $\omega_J$, and any such divisor $P$, there
is an actual source divisor $Q$ above $P$ such that
\begin{equation}
 m_Q=1,
 \qquad
 \frac{1+r_Q}{m_Q}
   =\inf_{T\mapsto P}\frac{1+r_T}{m_T},
 \qquad
 m_T=\ord_T(x^*P),\quad
 r_T=\frac{1}{p}\ord_T(\omega_J).
 \label{ds:threshold}
\end{equation}
Orders here are measured in the actual relative canonical bundle;
the infimum includes divisors on higher smooth models.  The same
statement holds after every finite parameter extension.

\begin{proposition}[Unramifiedness at moving divisors]
\label{ds:unramified}
The extension $\widetilde K/K$ in
Lemma~\ref{ds:cocycle} is unramified at every moving divisor.
After a finite extension $E/L_0$, each field component of
\[
 \widetilde K\otimes_K E(V')
\]
is likewise unramified at every moving divisor of $V'_E$.
\end{proposition}

\begin{proof}
We compare divisorial orders before and after splitting.  On the
constant model there is a unique minimizing divisor.  The actual
multiplicity-one minimizer in \eqref{ds:threshold} will force equality
of total and base inertia, and faithfulness will make both trivial.

\smallskip\noindent\textbf{Relative canonical orders.}\enspace
First work over $L_0$.  Put $L=K(J)$ and use
\eqref{ds:split-field} to identify
\[
 \widetilde L=L\widetilde K=\widetilde K(J_0).
\]
Since $J/K$ is geometrically integral, $L/K$ is regular.  Hence
$L$ and $\widetilde K$ are linearly disjoint over $K$, and
restriction identifies
\begin{equation}
 \operatorname{Gal}(\widetilde L/L)
      \xrightarrow{\sim}
 \operatorname{Gal}(\widetilde K/K)=G.
 \label{ds:regular-galois}
\end{equation}

Fix a moving divisor $P$ and a prolongation $P'$ to
$\widetilde K$.  Write $e=e(P'/P)$.  If a divisorial valuation
$Q$ of $L$ above $P$ prolongs to $Q'$ above $P'$, write
$e_Q=e(Q'/Q)$.  On smooth models at these generic divisor points,
the multiplicities and relative differential orders satisfy
\begin{equation}
 m_{Q'}=\frac{e_Qm_Q}{e},
 \qquad
 r_{Q'}=e_Qr_Q+e_Q-1-(e-1)m_{Q'}.
 \label{ds:ramification-orders}
\end{equation}
The first identity is the order of a base uniformizer.  For the
second, the total canonical differential acquires order $e_Q-1$
under the finite extension, whereas the base canonical
differential acquires order $e-1$ and is pulled back with
multiplicity $m_{Q'}$.  These are the tame different orders of
characteristic-zero divisorial DVRs.  They can also be computed
by differentiating uniformizers and completing them by separating
residue parameters.  Dividing the degree-$p$ calculation by $p$
gives the stated formula.

In particular,
\begin{equation}
 \frac{1+r_{Q'}}{m_{Q'}}
   =e\frac{1+r_Q}{m_Q}-(e-1).
 \label{ds:affine-threshold}
\end{equation}
This is an increasing affine transformation independent of $Q$.

\smallskip\noindent\textbf{The unique minimizing divisor.}\enspace
Let $R'$ be the valuation ring of $P'$ and let $\kappa'$ be its
residue field.  The valuation is trivial on $k$, because $P$ is
moving.  It is also trivial on the finite algebraic extension
$k'/k$: the valuation of a nonzero element algebraic over a
trivially valued field must be zero, as follows immediately from
its minimal polynomial.  Thus
\begin{equation}
 k'\subset R',\qquad k'\hookrightarrow\kappa'.
 \label{ds:constant-residue}
\end{equation}
Consider the smooth proper constant model
\[
 \mathcal J=J_0\times_{\operatorname{Spec}k}\operatorname{Spec}R'.
\]
Its special fiber $S'$ is smooth and geometrically integral over
$\kappa'$.  Choose the ordinary regular degree-$p$ generator
$\omega_0$ on $J_0$.  The pulled-back generator $\omega_J$ is a
scalar multiple $c\omega_0$ over $\widetilde K$, because the
ordinary degree-$p$ space has dimension one.  Put
$\ell=p^{-1}\ord_{P'}(c)$ and
$D_0=\operatorname{div}(\omega_0)\geq0$.

Every divisorial valuation of $\widetilde L$ above $P'$ has a
center on $\mathcal J$, by properness, and that center lies in
$S'$.  For such a valuation $T$, let $A_{\mathcal J}(T)$ denote its
log discrepancy over the smooth model $\mathcal J$.  Removing the
common scalar shift gives
\begin{equation}
 \frac{1+r_T}{m_T}
   =\ell+
     \frac{A_{\mathcal J}(T)+p^{-1}\ord_T(D_0)}{m_T},
 \qquad m_T=\ord_T(S').
 \label{ds:constant-threshold}
\end{equation}
The divisor $D_0$ is pulled back to the constant model in this
formula.  It has no special-fiber component.  Hence $S'$ itself
has $m_{S'}=1$ and ratio $\ell+1$.

For every other divisorial valuation centered in $S'$,
\[
 A_{\mathcal J}(T)>\ord_T(S')=m_T.
\]
This is the pure log terminality of a smooth divisor on a smooth
scheme.  In local parameters it follows by including the equation
of $S'$ and at least one further parameter vanishing at the proper
center: the log discrepancy is at least the sum of their positive
orders.  Since $\ord_T(D_0)\geq0$, Equation~\eqref{ds:constant-threshold}
proves that $S'$ is the unique
minimizer of the ratio above $P'$.

\smallskip\noindent\textbf{Triviality of inertia.}\enspace
Take the multiplicity-one minimizer $Q$ supplied by
\eqref{ds:threshold}.  It has a prolongation to $\widetilde L$;
by Galois transport we may arrange that the restriction of this
prolongation to $\widetilde K$ is the specified $P'$.  Indeed $G$
acts transitively on the prolongations of $P$, and its action on
$\widetilde L$ fixes $L$ by \eqref{ds:regular-galois}.
Conversely, every divisorial valuation above $P'$ restricts under
the finite extension to a divisorial valuation above $P$.
Equation \eqref{ds:affine-threshold} therefore shows that the
chosen prolongation of $Q$ attains the upstairs infimum.  The
uniqueness just proved identifies it with $S'$.  Consequently
$m_{Q'}=m_Q=1$, and \eqref{ds:ramification-orders} gives
\begin{equation}
 e_Q=e.
 \label{ds:equal-indices}
\end{equation}

Let $\mathcal I_{Q'}$ and $\mathcal I_{P'}$ be the total and base
inertia groups, viewed as subgroups of the common Galois group in
\eqref{ds:regular-galois}.  The embedding of base residue fields
in total residue fields gives
\[
 \mathcal I_{Q'}\subset\mathcal I_{P'}.
\]
All residue characteristics are zero, so the residue extensions
are separable and the orders of these inertia groups are $e_Q$
and $e$.  Equation \eqref{ds:equal-indices} forces equality of the
two inertia groups.  Thus every element of $\mathcal I_{P'}$ acts
trivially on the special-fiber function field $\kappa'(J_0)$.

For $g\in\mathcal I_{P'}$, coefficient transport is already
trivial on $\kappa'$.  Equation \eqref{ds:constant-residue} also
gives $g|_{k'}=1$.  Its induced action on $\kappa'(J_0)$ is
therefore the base extension of the constant map $z_g$.  Such an
action is faithful: a birational function-field automorphism over
$k'$ that becomes the identity after an injective extension of
fields was the identity to begin with.  Hence $z_g=1$.
The pair \eqref{ds:faithful-pairs} is faithful, so $g=1$.
This proves $\mathcal I_{P'}=1$, and hence unramifiedness at $P$.

\smallskip\noindent\textbf{Finite parameter extensions.}\enspace
Let $E/L_0$ be finite.  The base change of the finite
Galois splitting field is a finite product of fields, with a
transitive $G$-action; each component is Galois over $E(V')$ with
group its stabilizer in $G$.  The constant field $k'$ embeds in
every component.  The same constant maps and the restricted
faithful pairs define its semilinear action.  Geometric
integrality of $J$ and the multiplicity-one conclusion
\eqref{ds:threshold} both persist.  Repeating the preceding
argument for this component and its stabilizer proves the stated
unramifiedness after $E$.  In particular, no uniform choice of a
finite parameter extension is being assumed for all divisors.
\end{proof}

\subsection{Descent of the cover and the fiber}

\begin{proposition}[Descent of the whole fiber]
\label{ds:whole}
Let $\Lambda$ be an algebraic closure of $L_0=b(I)$.  There is a
finitely generated field $F_b/b$ containing $k=b(V')$ and an isomorphism
\begin{equation}
 \operatorname{Frac}(\Lambda\otimes_b F_b)
   \simeq
 \operatorname{Frac}(\Lambda\otimes_{L_0}K(J)).
 \label{ds:whole-field}
\end{equation}
The right-hand side is the function field of the whole geometric
generic fiber of the original fibration after identifying
$\Lambda$ with $\Omega$.  In particular,
\begin{equation}
 \Var(f)\leq\dim T_0.
 \label{ds:variation}
\end{equation}
\end{proposition}

\begin{proof}
\smallskip\noindent\textbf{A fixed open set carrying the cover.}\enspace
Extend the splitting algebra to $\Lambda(V')$:
\[
 \mathcal B_\Lambda
   =\widetilde K\otimes_K\Lambda(V').
\]
It is a finite product of fields and a Galois algebra with group
$G$.  We first descend its normalization over a suitable fixed
open subset of $V'$.

Every prime divisor of $V'_\Lambda$ is either the base extension
of a prime divisor of $V'$ or is moving. On an affine chart with
coordinate ring $A$, the map $A\to A\otimes_b\Lambda$ is faithfully
flat. If a height-one prime $\mathfrak q$ contracts to $\mathfrak p$,
flat going down gives $\operatorname{ht}\mathfrak p\leq1$.
A height-zero contraction is zero and means that the divisor is moving.
If $\operatorname{ht}\mathfrak p=1$, the corresponding integral
$b$-divisor is geometrically integral, since $b$ is algebraically
closed. Hence $\mathfrak p(A\otimes_b\Lambda)$ is a height-one
prime. Its inclusion in $\mathfrak q$ is equality. No finite-type
assumption on the field extension $\Lambda/b$ is used here.

A moving prime divisor of $V'_\Lambda$ descends to a moving prime
divisor after some finite extension $E/L_0$. Indeed, on the above
affine chart, choose finitely many generators of its prime ideal.
Their coefficients lie in such an $E$, since $\Lambda/L_0$ is
algebraic. The ideal generated by the same equations over $E$ is
prime by faithful-flat descent, and has height one by invariance of
codimension under field extension. Its contraction to $A$ is still
zero. Thus the descended divisor remains moving, and
Proposition~\ref{ds:unramified} applies to every field component at
this finite stage.

It follows that all divisorial branching of $\mathcal B_\Lambda$
lies on finitely many fixed divisors. The normalization is finite
because the varieties over a field are excellent; its branch locus
has only finitely many divisorial components.  Remove those divisors and
the singular locus of $V'$.  Also remove the fixed closed locus
over which the normalization in $k'/k$ is not finite \'etale.
This leaves a smooth nonempty open subset $O\subset V'$ over $b$.
The normalization of $O_\Lambda$ in $\mathcal B_\Lambda$ is
finite, normal, and unramified at all height-one points.  Purity
of the branch locus makes it finite \'etale.  Here the base is
regular, the source is normal, the morphism is finite and
generically separable, and all divisorial ramification has been
removed, which are the hypotheses of Zariski--Nagata purity;
see \cite[Expos\'e~X, Th\'eor\`eme~3.1]{SGA1}.
Denote this finite \'etale cover by
\[
 \mathcal C_\Lambda\longrightarrow O_\Lambda.
\]
Its generic $G$-action extends by normalization.  It is a
$G$-torsor, possibly disconnected, since the generic Galois
algebra is a $G$-torsor and all these covers are finite \'etale
over the connected base.

Let $\mathcal C_{1,\Lambda}$ be a connected component, and let
$G_1\subset G$ be its stabilizer.  Then
$\mathcal C_{1,\Lambda}\to O_\Lambda$ is a connected finite
\'etale $G_1$-torsor.  Let $\mathcal T\to O$ be the finite
\'etale cover obtained by normalization in $k'/k$.  The inclusion
$k'\subset\mathcal B_\Lambda$ gives a morphism
\[
 \mathcal C_{1,\Lambda}\longrightarrow\mathcal T_\Lambda
\]
over $O_\Lambda$, compatible with the action of $G_1$ through its
restriction to $k'$.

\smallskip\noindent\textbf{Descent with the coefficient field.}\enspace
Finite \'etale covers are invariant under an extension of
algebraically closed fields in characteristic zero.  In the
needed form, base extension from $b$ to $\Lambda$ gives an
equivalence between the categories of finite \'etale covers of
the connected normal separated finite-type scheme $O$ and of
$O_\Lambda$.  This follows from
\cite[Theorem~1.1]{Landesman24}, since every finite cover is tame in
characteristic zero; see also \cite[Remark~1.5]{Landesman24} and
\cite[Expos\'e~XIII, Proposition~4.6]{SGA1}.  Full faithfulness
descends the morphisms as well as the covers.

Apply this equivalence simultaneously to
$\mathcal C_{1,\Lambda}$, its $G_1$-action, and its map to the
fixed cover $\mathcal T_\Lambda$.  We obtain a connected finite
\'etale $G_1$-torsor $\mathcal C_1\to O$ and a compatible map
$\mathcal C_1\to\mathcal T$.  Let $E_1$ be the function field of
$\mathcal C_1$.  Then
\begin{equation}
 E_1/k\ \text{is finite Galois with group }G_1,
 \qquad k'\subset E_1,
 \label{ds:cover-field}
\end{equation}
and the chosen component of $\mathcal B_\Lambda$ is
$\operatorname{Frac}(\Lambda\otimes_bE_1)$.
The embedding of $k'$ and the restriction of the $G_1$-action to
it are part of this descended diagram.

\smallskip\noindent\textbf{The invariant field.}\enspace
For $g\in G_1$, the birational map $z_g$ is defined over $k'$,
and $k'\subset E_1$.  Hence the semilinear action
$z_gs_g$ descends to the field $E_1(J_0)$.  Its multiplication
law follows from the cocycle law, because the action on $k'$ in
\eqref{ds:cover-field} is the descended one.  Define
\begin{equation}
 F_b=E_1(J_0)^{G_1}.
 \label{ds:descended-field}
\end{equation}
This is a finitely generated field over $b$.  Indeed
$E_1(J_0)$ is finitely generated over $b$ and finite over its
invariant field.  Also
\begin{equation}
 k=b(V')\subset F_b,
 \label{ds:retained-coordinates}
\end{equation}
since both the coefficient action of $G_1$ and the constant
birational maps fix $k$.  This inclusion retains the coordinates
of the entire $x$-base in the descended field.

The coefficient field and its action have both descended. We now
verify \eqref{ds:whole-field}, including its retained $k$-coordinates.  Invariants on a finite product
of fields permuted transitively by a finite group are identified
with invariants on a chosen factor under its stabilizer: the
other coordinates of an invariant element are its translates.
Thus after extension to $\Lambda$, Equation~\eqref{ds:split-field}
may be computed on the chosen component
with group $G_1$.

Invariants also commute with extension from $b$ to $\Lambda$.
For a ring with a finite $G_1$-action this follows by applying the
averaging projector $|G_1|^{-1}\sum_{g\in G_1}g$.  Passage to
fraction fields causes no change: if an invariant fraction has
denominator $a$, multiplication by the invariant denominator
$\prod_{g\in G_1}g(a)$ puts its numerator in the invariant ring.
These arguments apply to the present tensor products, which are
domains on the chosen component; the finitely generated fields
over the algebraically closed field $b$ are geometrically
integral.  Consequently
\[
 \operatorname{Frac}(\Lambda\otimes_b F_b)
  \simeq
 \left(
   \operatorname{Frac}(\Lambda\otimes_b E_1(J_0))
 \right)^{G_1}
  \simeq
 \operatorname{Frac}(\Lambda\otimes_{L_0}K(J)).
\]

\smallskip\noindent\textbf{Identification with the original fiber.}\enspace
Proposition~\ref{it:dimensions} gives $\C(S)=\C(Y)$ in the
original field tower
\[
 \C(Y)\subset\C(S)\subset\C(W)\subset\C(X).
\]
In particular, $\C(T_0)\subset\C(Y)$, with its prescribed
algebraic closure $b\subset\Omega$.  Before the finite parameter
changes the generic parameter field is
\[
 b(I)=b\,\C(Y)\subset\Omega.
\]
It is algebraic over $\C(Y)$.  Each subsequent finite parameter
extension can be embedded in $\Omega$, and an algebraic closure
of the resulting $L_0$ is still $\Omega$.  We may therefore take
$\Lambda=\Omega$.

By Proposition~\ref{mk:product}, after this parameter change the
function field of the generic $x$-base is $L_0(V')=K$.
The field $K(J)$ is its total field in the original fibration.
Birational model changes preserve these fields, so
\[
 \operatorname{Frac}(\Omega\otimes_{L_0}K(J))
   \simeq
 \operatorname{Frac}
       (\Omega\otimes_{\C(Y)}\C(X)).
\]
The right-hand side is the function field of the whole smooth
geometric generic fiber of $f$, not just that of $J$ over an
algebraic closure of $k$.

Finally, take a smooth projective $b$-model of $F_b$, by a
projective compactification and resolution of singularities in
characteristic zero.  Equation \eqref{ds:whole-field} says that
its base extension to $\Omega$ is birational to that whole fiber.
Since
\[
 \trdeg_{\C}b=\dim T_0,
\]
the definition of birational variation gives
\eqref{ds:variation}.
\end{proof}

\subsection{The variation inequality}

\begin{proof}[Proof of Theorem~\ref{thm:main}]
If $\kappa(F)=-\infty$,
the assertion holds by the stated convention.  Otherwise set
$d=\kappa(F)\geq0$.  The compactification and logarithmic
subadditivity argument in Lemma~\ref{it:compactification},
specifically \eqref{it:base-bound}, gives
\[
 \kbar(U)\geq d+\kbar(V).
\]
Equations \eqref{it:iitaka-dimension} and
\eqref{it:dimension-equations} identify the dimension of the
absolute logarithmic Iitaka base and give
\[
 \kbar(U)=\dim Z=d+\dim T_0.
\]
Combining this identity with \eqref{ds:variation} gives
$\kbar(U)\geq d+\Var(f)$.  Taking the larger of the two lower
bounds proves
\[
 \boxed{\displaystyle
 \kbar(U)\geq\kappa(F)
                  +\max\{\kbar(V),\Var(f)\}.}
\]

The dimension-zero cases fit the same field argument.  If $d=0$,
then $V'$ is a point over $b$, so $k=b$ and the geometric
constancy already gives the required descent.  If $I$ is a
point, there are no parameter directions and $\Lambda=b$.
If $J_0$ is a point, its birational group is trivial and the
descended field retains precisely the fixed $V'$ coordinates.
For a point base or a zero-dimensional fiber the same conventions
give the asserted inequality.  Throughout, the Kodaira dimension
used for the original compact generic fiber is its ordinary
Kodaira dimension, and every auxiliary marking has been removed
from the final invariant field.
\end{proof}

When $U$ and $V$ are projective, their logarithmic Kodaira dimensions
are their ordinary Kodaira dimensions.  The same inequality therefore
gives the ordinary Iitaka--Viehweg $C^+$ statement.
